\documentclass[10pt,a4paper]{amsart}
\usepackage[margin=19mm]{geometry}
\usepackage{amsmath,amssymb,mathtools}
\usepackage[scr=rsfs,cal=euler]{mathalfa}
\usepackage{xfrac}
\usepackage[unicode,hidelinks,bookmarks=false]{hyperref}
\newcommand{\ie}{i.e.\ }
\newcommand{\cf}{cf.\ }
\newcommand{\resp}{respectively\ }
\newcommand{\Subsection}{\subsection}
\providecommand{\nobreakdash}{\nobreak\hbox{-}}
\setlength{\emergencystretch}{2em}
\multlinegap0pt
\usepackage{bm}
\usepackage{bbm}
\usepackage{dsfont}
\usepackage{xspace}
\usepackage{cancel}
\usepackage{mathtools}
\usepackage{amsthm}
\makeatletter
\@ifundefined{theorem}{\newtheorem{theorem}{Theorem}}{}
\@ifundefined{lemma}{\newtheorem{lemma}[theorem]{Lemma}}{}
\makeatother
\DeclareMathOperator{\CC}{\mathtt{C}}
\newcommand{\CS}[1][]{\ensuremath{\mathfrak{C}_{#1}}}
\title[NOC NOC, who's there? Clustering systems of tree-child and n]{NOC NOC, who's there? Clustering systems of tree-child and normal networks}
\author{Anna Lindeberg, Marc Hellmuth}
\date{Source: arXiv:2609.13336v1 (CC BY 4.0)}
\begin{document}
\maketitle

\noindent\emph{Source and licence.} NOC NOC, who's there? Clustering systems of tree-child and normal networks. Version arXiv:2609.13336v1 (CC BY 4.0), distributed under the Creative Commons Attribution 4.0 International licence, as recorded in the arXiv licence metadata for this exact version and shown on the abstract page https://arxiv.org/abs/2609.13336v1. Full rights notice: licenses/arxiv-2609.13336-CC-BY-4.0.txt.\par\smallskip

\noindent\textit{Editorial scope.} Counted unit: the Lemma whose source label is \texttt{lem:visible-semiregular} in the article (source0.tex lines 268--272, source label \texttt{lem:visible-semiregular}), delivered together with the article's own complete original proof of it (source0.tex lines 273--300, one \texttt{proof} environment of 28 lines). No mathematics is written, repaired or re-derived: statement and proof are the article's own text line for line. Required context retained uncounted: none. Every definition, display and label of those lines is retained unchanged. Omitted from this package and not redistributed: 1--267, 301--674. Nothing inside a retained excerpt is omitted or altered: every retained body is delivered exactly as the article's lines are written, so no omission receipt of any kind accompanies this package. Citations: the retained excerpts carry no citation command at all, so no bibliography entry of the article is transcribed and no reference is redistributed. Reference adaptation: none was needed and none was made; every reference inside the delivered unit resolves inside it, so no unresolved reference or citation is produced. Printed slips kept literal: no printed word, symbol, hypothesis or number of the counted statement or of its proof was altered. Layout and class: the article's own class is \texttt{elsarticle} with packages including geometry, hyperref, booktabs, amsmath, amsfonts, amsthm, amssymb, mathtools, mathrsfs, xcolor; the wrapper is the lane's pinned \texttt{amsart} editorial wrapper at 10pt on A4, which restates the theorem-like environments the retained text needs and adds the article's own macro definitions that the retained text needs (\texttt{\textbackslash CC}, \texttt{\textbackslash CS}), with extra wrapper packages (bm, bbm, dsfont, xspace, cancel, mathtools). The delivered numbering is the unit's own. Assets: no retained span contains a figure environment, a graphics-inclusion command, a PDF-page inclusion, a table or a TikZ picture; no image file is copied into this package, the package declares no asset dependency and no image file is redistributed with it. Licence: version arXiv:2609.13336v1 of this article is distributed under the Creative Commons Attribution 4.0 International licence, as recorded in the arXiv licence metadata for this exact version and shown on the abstract page https://arxiv.org/abs/2609.13336v1. A full rights notice is delivered at licenses/arxiv-2609.13336-CC-BY-4.0.txt. Characters: every retained excerpt is pure ASCII, so the delivered engine, XeLaTeX with its default Latin Modern text font, reports no missing character.
\begin{lemma}\label{lem:visible-semiregular}
    Let $N$ be a network and let $v\in V(N)$. If $v$ is visible, then
    $\CC_N(v)$ is inclusion-visible in $\CS[N]$. If $N$ is semi-regular, then $v$ is
    visible if and only if $\CC_N(v)$ is inclusion-visible in $\CS[N]$.
\end{lemma}

\begin{proof}
	Let $N$ be a network with root $\rho$, let $v\in V(N)$ and put $C=\CC_N(v)$. First assume that
	$v$ is visible.
	Then, there exists some $x\in C$ such that every $\rho x$-path in $N$ contains $v$. Let
	$D\in\CS[N]$ be any cluster with $x\in D$. Then, $D=\CC_N(w)$ for some $w\in V(N)$.
	Since $x\in D$, there is a
	directed path from $w$ to $x$, and hence a $\rho x$-path $P'$ containing $w$, obtained by
	concatenating any $\rho w$-path with any $wx$-path.
	Since every $\rho x$-path in $N$ contains $v$, it holds that $v$ is a vertex of $P'$, and
	hence, $v\preceq_N w$ or $w\preceq_N v$. In either case, we have that either $C\subseteq D$ or
	$D\subseteq C$. Hence, $C$ is inclusion-visible in $\CS[N]$.
	
	Now, assume that $N$ is semi-regular. By the first part, it remains to prove the
	\emph{if-direction} of the second statement.
 	Hence, suppose that $\CC_N(v)$ is inclusion-visible in $\CS[N]$. Then, there exists some
	inclusion-visible witness $x$ for $C$. Thus,
	for all $D\in\CS[N]$ with $x\in D$, we have $D\subseteq C$ or $C\subseteq D$.
	Let $P=v_0v_1\ldots v_k$ be an arbitrary $\rho x$-path in $N$, where $v_0=\rho$ and $v_k=x$.
	Since $x\preceq_N v_i$, we have  $x\in \CC_N(v_i)$ for all $i\in\{0,\ldots,k\}$. 
	By the choice of $x$, it follows that $\CC_N(v_i)\subseteq \CC_N(v)$ or
	$\CC_N(v)\subseteq \CC_N(v_i)$ for all $i\in\{0,\ldots,k\}$. Since $N$ satisfies 
	(PCC), the vertices $v$ and $v_i$ are therefore $\preceq_N$-comparable for all
	$i\in\{0,\ldots,k\}$. The latter together with the fact that $x=v_k\preceq_N v\preceq_N v_0=\rho$
	implies that there is some $j\in\{0,\ldots,k-1\}$ such that $v_{j+1}\preceq_N v\preceq_N v_j$.
	Since $v_{j+1}$ and $v_j$ are adjacent in $N$ and $N$ has no shortcuts, we conclude that
	$v=v_{j+1}$ or $v=v_j$. In either case, $v$ is a vertex of $P$. Since $P$ was an
	arbitrary $\rho x$-path in $N$, we conclude that $v$ is a visible vertex.
\end{proof}

\end{document}
